\documentclass[11pt,spanish]{article}

% =========================================================
% PLANTILLA BASE PARA TAREAS ACADÉMICAS / INFORMES TÉCNICOS
% =========================================================

% --- Idioma y codificación para Tectonic / XeLaTeX ---
\usepackage[spanish,es-tabla,es-noshorthands]{babel}
\usepackage{fontspec}

% --- Márgenes / papel ---
\usepackage{geometry}
\geometry{a4paper, margin=2.7cm}

% --- Matemáticas ---
\usepackage{amsmath,amssymb}

% --- Tablas, enlaces y elementos visuales ---
\usepackage{array}
\usepackage{tabularx}
\usepackage{booktabs}
\usepackage{longtable}
\usepackage{graphicx}
\usepackage{float}
\usepackage{xcolor}
\usepackage{tikz}
\usepackage{pgfplots}
\pgfplotsset{compat=1.16}
\usetikzlibrary{positioning,shapes.geometric,arrows.meta,calc}
\usepackage{enumitem}
\usepackage{ragged2e}
\usepackage{caption}
\captionsetup{skip=6pt}
\captionsetup[table]{position=top}
\captionsetup[figure]{position=bottom}
\usepackage{xurl}
\usepackage[hidelinks]{hyperref}

% --- Configuración de párrafo en bloque ---
\makeatletter
\renewcommand\paragraph{\@startsection{paragraph}{4}{\z@}%
  {2.2ex \@plus1ex \@minus.2ex}%
  {0.8ex}%
  {\normalfont\normalsize\bfseries}}
\makeatother

% --- Listados de código con formato elegante ---
\usepackage{listings}

\definecolor{codebg}{rgb}{0.97,0.97,0.98}
\definecolor{codeframe}{rgb}{0.85,0.85,0.88}
\definecolor{commentgreen}{rgb}{0.13,0.54,0.13}
\definecolor{keywordblue}{rgb}{0.0,0.2,0.7}
\definecolor{stringpurple}{rgb}{0.6,0.1,0.4}

% --- Colores para esquemas y diagramas TikZ ---
\definecolor{softblue}{rgb}{0.88, 0.93, 1.0}
\definecolor{softgreen}{rgb}{0.88, 0.96, 0.88}
\definecolor{softorange}{rgb}{1.0, 0.92, 0.82}
\definecolor{softpurple}{rgb}{0.93, 0.88, 0.98}
\definecolor{softteal}{rgb}{0.85, 0.95, 0.95}
\definecolor{softyellow}{rgb}{1.0, 0.95, 0.80}

\lstdefinestyle{customcode}{
    backgroundcolor=\color{codebg},
    basicstyle=\ttfamily\footnotesize,
    breakatwhitespace=false,
    breaklines=true,
    captionpos=b,
    keepspaces=true,
    numbers=none,
    showspaces=false,
    showstringspaces=false,
    showtabs=false,
    tabsize=4,
    frame=single,
    rulecolor=\color{codeframe},
    commentstyle=\color{commentgreen}\itshape,
    keywordstyle=\color{keywordblue}\bfseries,
    stringstyle=\color{stringpurple},
    xleftmargin=0.5em,
    xrightmargin=0.5em
}

\lstdefinestyle{customcodenum}{
    backgroundcolor=\color{codebg},
    basicstyle=\ttfamily\footnotesize,
    breakatwhitespace=false,
    breaklines=true,
    captionpos=b,
    keepspaces=true,
    numbers=left,
    numbersep=6pt,
    numberstyle=\tiny\color{gray},
    showspaces=false,
    showstringspaces=false,
    showtabs=false,
    tabsize=4,
    frame=single,
    rulecolor=\color{codeframe},
    commentstyle=\color{commentgreen}\itshape,
    keywordstyle=\color{keywordblue}\bfseries,
    stringstyle=\color{stringpurple},
    xleftmargin=1.5em,
    xrightmargin=0.5em
}

\lstset{style=customcode}

% --- Encabezado y pie de página: solo número de página en el pie ---
\usepackage{fancyhdr}
\setlength{\headheight}{15pt}
\setlength{\footskip}{25pt}

\pagestyle{fancy}
\fancyhf{}
\renewcommand{\headrulewidth}{0pt}
\renewcommand{\footrulewidth}{0pt}
\fancyfoot[C]{\thepage}

\fancypagestyle{plain}{%
    \fancyhf{}%
    \renewcommand{\headrulewidth}{0pt}%
    \renewcommand{\footrulewidth}{0pt}%
    \fancyfoot[C]{\thepage}%
}

% --- \input de filas generadas dentro de tabular (el \input de LaTeX añade ganchos que rompen \bottomrule) ---
\makeatletter\newcommand{\filas}[1]{\@@input #1 }\makeatother

% --- Interlineado ---
\linespread{1.2}
\sloppy

% =========================================================
% DATOS DEL DOCUMENTO
% =========================================================
\newcommand{\Institucion}{Instituto Tecnológico de Costa Rica}
\newcommand{\Campus}{Escuela de Ingeniería en Computación}
\newcommand{\DocumentoTitulo}{Progra01}
\newcommand{\DocumentoSubtitulo}{Creación del índice de búsqueda}
\newcommand{\NombreBuscador}{Encuadre}
\newcommand{\Curso}{IC-8060 Recuperación de Información Textual}
\newcommand{\Profesor}{Aurelio Sanabria}
\newcommand{\FechaDocumento}{Noviembre de 2026}
\newcommand{\PeriodoAcademico}{II Semestre}
\newcommand{\AnioDocumento}{2026}

% Autores con nombre a la izquierda y carné a la derecha
\newcommand{\AutoresPortada}{%
\begin{minipage}{0.85\textwidth}
\begin{tabular*}{\linewidth}{@{\extracolsep{\fill}}lr@{}}
Santiago Espinoza Rendon & 2024156530 \\
Fabián Araya Ortega & 2019159110 \\
Mauricio González Prendas & 2024143009 \\
% José Fabián Zumbado Ruiz & <carné> \\  % ← completar el carné y descomentar
\end{tabular*}
\end{minipage}%
}

% Metadatos PDF
\title{\DocumentoTitulo}
\author{Santiago Espinoza Rendon, Fabián Araya Ortega, Mauricio González Prendas}
\date{\FechaDocumento}

% =========================================================
% CONFIGURACIÓN DE TABLAS Y UTILIDADES
% =========================================================
\newcolumntype{Y}{>{\RaggedRight\arraybackslash}X}
\newcolumntype{P}[1]{>{\RaggedRight\arraybackslash}p{#1}}
\newcommand{\offset}[1]{{\ttfamily\scriptsize\color{blue!80!black}#1}}
\renewcommand{\arraystretch}{1.22}
\setlength{\LTpre}{0.5em}
\setlength{\LTpost}{0.5em}

% =========================================================
% PORTADA
% =========================================================
\newcommand{\MakeCover}{%
\begin{titlepage}
\thispagestyle{empty}
\centering

{\bfseries \huge \Institucion \par}
\vspace{0.25cm}
{\LARGE \Campus \par}
\vspace{0.15cm}
{\LARGE \Curso \par}
\vspace{0.5cm}
\includegraphics[width=0.22\linewidth]{imgs/logo.png}

\vfill

{\huge \bfseries \DocumentoTitulo: \DocumentoSubtitulo \par}
\vspace{0.4cm}
{\Large \textit{\NombreBuscador} \par}

\vfill

{\bfseries \LARGE Profesor: \par}
\vspace{0.2cm}
{\Large \Profesor \par}

\vfill

{\bfseries \LARGE Estudiantes: \par}
\vspace{0.3cm}
{\Large \AutoresPortada \par}

\vfill

{\bfseries\LARGE Fecha: \par}
\vspace{0.2cm}
{\Large \FechaDocumento \par}

\vfill

{\LARGE \PeriodoAcademico \par}
{\LARGE \AnioDocumento \par}

\end{titlepage}%
}

% =========================================================
% DOCUMENTO PRINCIPAL
% =========================================================
\begin{document}
\hypersetup{pageanchor=false}

% =========================
% Portada
% =========================
\MakeCover

% =========================
% Índice
% =========================
\pagenumbering{roman}
\pagestyle{plain}
\tableofcontents
\clearpage
\pagenumbering{arabic}
\hypersetup{pageanchor=true}
\pagestyle{fancy}

% =========================================================
% SECCIÓN 1: INTRODUCCIÓN
% =========================================================
\section{Introducción}

Este documento describe cómo se normaliza el texto antes de indexarlo y qué efecto tiene esa decisión sobre el índice. Una consulta solo encuentra un documento si ambos se reducen a los mismos términos; por eso la normalización define qué significa que dos palabras sean ``la misma'' y condiciona tanto el tamaño del índice como la calidad de los resultados~\cite{manning2008}.

Se definieron tres configuraciones de normalización (mínima, intermedia y agresiva) y se construyó con cada una un índice sobre la misma muestra de la colección. La sección~\ref{sec:politicas} explica cada política y su justificación, la sección~\ref{sec:estadisticas} presenta las estadísticas medidas y la sección~\ref{sec:estructura-binaria} describe el formato binario del archivo. Las listas de palabras vacías, sinónimos y tesauro viven en \texttt{recursos/} y las leen por igual el indexador propio y el de Lucene; los casos de prueba compartidos están en \texttt{casos\_prueba/normalizacion.json}.

% =========================================================
% SECCIÓN 2: POLÍTICAS DE NORMALIZACIÓN
% =========================================================
\section{Políticas de normalización}
\label{sec:politicas}

Cada política responde a tres preguntas: qué se hace, por qué, y dónde se aplica en cada indexador. La consulta recibe exactamente la misma normalización que el documento; si no fuera así, un término buscado nunca coincidiría con el indexado. El buscador reconstruye la normalización del índice a partir de las banderas guardadas en el encabezado del archivo, de modo que una consulta contra la política mínima no se reduce como la intermedia.

\subsection{Tokenización y limpieza previa (P1)}
\textbf{Descripción.} El texto se pasa a la forma Unicode NFKC, se unifican las comillas y rayas tipográficas, y se parte en tokens con la expresión regular \texttt{<[A-Z0-9\_]+>|[\textasciicircum\textbackslash W\_]+(?:[-'\_][\textasciicircum\textbackslash W\_]+)*}. Una palabra con guion o apóstrofe interno (\textit{sci-fi}, \textit{spider-man}, \textit{o'brien}) es un solo token, y los marcadores sintéticos como \texttt{<NUM>} se conservan enteros.

\textbf{Justificación.} NFKC convierte ligaduras y caracteres de ancho completo a su forma ordinaria, de modo que cadenas que se ven igual quedan con la misma representación binaria~\cite{unicode15}. Partir \textit{sci-fi} en dos tokens perdería el término que el usuario escribe, porque ninguna de las dos mitades significa lo mismo que la palabra completa.

\textbf{Dónde.} Propio: \texttt{normalizador.py}, paso de limpieza y \texttt{RE\_TOKEN}. Lucene: \texttt{PatternReplaceCharFilter} para las comillas y un \texttt{PatternTokenizer} con la misma expresión (la sintaxis de Java la acepta sin cambios).

\subsection{Minúsculas (P2)}
\textbf{Descripción.} Todo token alfanumérico pasa a minúsculas; los marcadores \texttt{<...>} no se tocan.

\textbf{Justificación.} \textit{Cinema}, \textit{CINEMA} y \textit{cinema} deben ser una sola entrada del diccionario; distinguirlas solo dispersaría los postings~\cite{manning2008}. El costo es que se pierde la diferencia entre un nombre propio y una palabra común (\textit{Hook} frente a \textit{hook}), pérdida que se asume porque la colección es de texto libre con capitalización irregular.

\textbf{Dónde.} Propio: \texttt{str.lower}. Lucene: \texttt{LowerCaseFilter}.

\subsection{Palabras vacías (P3)}
\textbf{Descripción.} Se descartan los tokens que aparecen en \texttt{recursos/stopwords\_es.txt} o \texttt{stopwords\_en.txt} (una palabra por línea; \texttt{\#} inicia un comentario). Del inglés también se quitan las contracciones (\textit{'s, 're, 've, 'll, 'd, 'm}) y las formas \textit{n't}, de modo que \textit{it's} produce \textit{it} y este se descarta como vacía. Las dos listas se aplican a todo documento, sin importar su idioma, porque la colección mezcla español e inglés dentro de un mismo texto.

\textbf{Justificación.} Las palabras de altísima frecuencia ocupan las listas de postings más largas y casi no discriminan entre documentos; con ponderación $tf\text{-}idf$ su peso ya es cercano a cero~\cite{salton1988}. Quitarlas reduce el índice a costa de impedir consultas que sean solo de palabras vacías, lo que en este dominio no tiene valor. Se evitaron entradas ambiguas con el dominio: por ejemplo, no se incluyó \textit{star} porque forma parte de \textit{Star Wars}.

\textbf{Dónde.} Propio: \texttt{cargar\_lista\_palabras}. Lucene: \texttt{WordlistLoader.getWordSet(reader, "\#")} con \texttt{StopFilter}~\cite{lucene-wordlist}.

\subsection{Reducción morfológica (P4)}
\textbf{Descripción.} Se aplica el algoritmo Snowball del idioma del documento. El idioma se decide contando, entre los primeros 300 tokens, cuántos están en cada lista de palabras vacías; gana el inglés solo si supera estrictamente al español, y un empate o la ausencia de votos da español. Un documento usa un único stemmer.

\textbf{Justificación.} La reducción unifica formas flexivas (\textit{movies}, \textit{movie}) y mejora la cobertura, pero es una heurística: Porter muestra que se trata de quitar sufijos con reglas, no de hallar la forma de diccionario~\cite{porter1980,porter2001}. Por eso en este documento se habla de \textbf{stemming}, no de lematización. Los resultados empíricos de Harman muestran que reducir sufijos no mejora de forma consistente la recuperación en colecciones pequeñas~\cite{harman1991}, así que la política mínima se conserva como línea base para medir el efecto real sobre esta colección.

\textbf{Dónde.} Propio: \texttt{snowballstemmer} con memoria de las raíces ya calculadas. Lucene: \texttt{SnowballFilter} con el mismo idioma; otros stemmers, como \texttt{PorterStemFilter}, no son equivalentes y no deben usarse.

\subsection{Tildes y la letra eñe (P5)}
\textbf{Descripción.} En la política agresiva se pliegan las tildes y diéresis (\textit{película} $\to$ \textit{pelicula}), pero la \textit{ñ} se conserva: \textit{año} y \textit{ano} siguen siendo distintas.

\textbf{Justificación.} Quien busca sin tildes debe encontrar texto con tildes. En cambio \textit{año}/\textit{ano} son palabras distintas y colapsarlas crea ruido real. Plegar solo las tildes de las vocales es la excepción que evita ese error.

\textbf{Dónde.} Propio: \texttt{remover\_diacriticos}. Lucene: un \texttt{MappingCharFilter} propio; \texttt{ASCIIFoldingFilter} no sirve tal cual porque también convierte \textit{ñ} en \textit{n}.

\subsection{Fechas (P6)}
\textbf{Descripción.} En la política agresiva, cada fecha completa (\texttt{25/12/1977}, \texttt{1977-12-25}, \textit{25 de diciembre de 1977}, \textit{December 25, 1977}) se sustituye por \texttt{<FECHA>}. Un año suelto entre 1888 y 2099 se sustituye por \texttt{<ANIO\_aaaa>}, de modo que el año de estreno sigue siendo buscable.

\textbf{Justificación.} Una fecha completa casi nunca se busca tal cual, pero ocupa un término distinto por cada día del calendario; el año sí es una consulta frecuente (\textit{1994}). El rango 1888--2099 acota los años plausibles para una obra audiovisual y evita convertir en año cualquier número de cuatro cifras.

\textbf{Dónde.} Propio: \texttt{RE\_FECHA\_TEXTUAL} y las expresiones numéricas. Lucene: \texttt{PatternReplaceCharFilter} con las mismas expresiones.

\subsection{Números (P7)}
\textbf{Descripción.} Un número (\texttt{\textbackslash d+([.,]\textbackslash d+)?}) se puede conservar, descartar o reemplazar por \texttt{<NUM>}; la política agresiva usa \texttt{<NUM>}. Los años y las fechas ya se trataron en P6 y no llegan a este paso.

\textbf{Justificación.} Los números de la colección (duraciones, puntajes, minutos, versiones) casi no sirven como consulta, pero llenan el vocabulario con términos únicos. Con \texttt{<NUM>} se conserva la información de que hubo un número sin pagar un término por cada valor.

\textbf{Dónde.} Propio: \texttt{RE\_NUMERO}. Lucene: \texttt{PatternReplaceCharFilter}.

\subsection{Sinónimos y tesauro (P8)}
\textbf{Descripción.} Dos archivos con el mismo formato (\texttt{variante, variante => descriptor}) llevan variantes a un solo descriptor. \texttt{sinonimos.txt} cubre equivalencias de obras, formatos y técnica (\textit{película, movie, cine, filme} $\to$ \texttt{film}; \textit{guion, script} $\to$ \texttt{screenplay}; \textit{soundtrack, ost} $\to$ \texttt{banda\_sonora}). \texttt{tesauro.txt} cubre oficios y créditos (\textit{director, cineasta, filmmaker} $\to$ \texttt{realizador}; \textit{actor, actriz, intérprete} $\to$ \texttt{interprete}). Solo se mapean tokens sueltos y, si una variante está en ambos archivos, gana el tesauro. Los descriptores se reducen con el stemmer de su propio idioma, para que \texttt{film} y \texttt{realizador} den el mismo término en cualquier documento.

\textbf{Justificación.} Un tesauro es un vocabulario controlado: un descriptor preferido y términos no preferidos que remiten a él~\cite{iso25964}. Aplicarlo al indexar y al consultar permite que una consulta en español encuentre un documento en inglés. Los resultados de expansión léxica con WordNet fueron pobres cuando se aplicó a ciegas~\cite{voorhees1994}; por eso las listas son cortas, escritas a mano y solo incluyen variantes que significan lo mismo en el dominio. Se excluyeron las ambiguas: \textit{star} (Star Wars), \textit{score} (puntaje), \textit{short}, \textit{feature}, \textit{editor}, \textit{composer} y \textit{protagonista}.

\textbf{Dónde.} Propio: \texttt{cargar\_sinonimos}. Lucene: \texttt{SolrSynonymParser} para ambos archivos~\cite{lucene-solrsyn} con \texttt{SynonymGraphFilter}, usado como reemplazo (\texttt{expand=false}).

\subsection{Nombres propios (P9)}
\textbf{Descripción.} No se hace reconocimiento de entidades. Un nombre propio recibe el mismo trato que cualquier otro token: minúsculas, tildes (según la política) y, con stemming, reducción. \textit{George Lucas} produce \texttt{georg luca}.

\textbf{Justificación.} El reconocimiento de entidades nombradas es una tarea con sus propios modelos y su propia evaluación~\cite{nadeau2007}, que excede el alcance del proyecto. Como la consulta pasa por la misma reducción, \textit{George Lucas} sigue encontrando \textit{George Lucas}; lo que se pierde es distinguir entre un apellido y una palabra común, y que el stemmer puede unir nombres distintos con la misma raíz.

\subsection{Las tres configuraciones}

La Tabla~\ref{tab:politicas} resume qué activa cada configuración. Cada una es un archivo \texttt{.conf} en \texttt{configuraciones/}.

\begin{table}[htbp]
  \centering
  \footnotesize
  \caption{Pasos activos en cada configuración de normalización.}
  \label{tab:politicas}
  \begin{tabular*}{\linewidth}{@{\extracolsep{\fill}}lccc@{}}
    \toprule
    \textbf{Paso} & \textbf{Mínima} & \textbf{Intermedia} & \textbf{Agresiva} \\
    \midrule
    Limpieza Unicode y tokenización (P1) & \checkmark & \checkmark & \checkmark \\
    Minúsculas (P2) & \checkmark & \checkmark & \checkmark \\
    Contracciones y palabras vacías (P3) & -- & \checkmark & \checkmark \\
    Stemming Snowball (P4) & -- & \checkmark & \checkmark \\
    Plegado de tildes, conserva \textit{ñ} (P5) & -- & -- & \checkmark \\
    Fechas y años (P6) & -- & -- & \checkmark \\
    Números a \texttt{<NUM>} (P7) & -- & -- & \checkmark \\
    Sinónimos y tesauro (P8) & -- & -- & \checkmark \\
    \midrule
    Archivo & \texttt{minimo.conf} & \texttt{estandar.conf} & \texttt{avanzado.conf} \\
    \bottomrule
  \end{tabular*}
\end{table}

\textbf{Mínima.} Ventaja: es la línea base, no pierde ninguna palabra y cualquier consulta literal funciona. Desventaja: el vocabulario es el más grande y las palabras vacías llenan las listas de postings.

\textbf{Intermedia.} Ventaja: reduce el vocabulario y el archivo sin alterar el significado de las palabras. Desventaja: el stemming une palabras que no deberían unirse y una consulta sin tildes no encuentra texto con tildes.

\textbf{Agresiva.} Ventaja: mayor reducción y tolera consultas sin tildes y en otro idioma gracias al tesauro. Desventaja: pierde la fecha y los valores numéricos exactos (incluido el año cuando forma parte de una fecha completa), y cada regla del tesauro es una decisión editorial que puede equivocarse.

\subsection{Limitaciones}
\begin{itemize}
  \item Un documento usa un solo stemmer; un guion con diálogos en los dos idiomas se reduce con el del idioma dominante.
  \item Las reglas de sinónimos son de un solo token: no hay frases como \textit{banda sonora} $\to$ \texttt{banda\_sonora}, porque el tokenizador ya separó las palabras.
  \item Los nombres propios no se detectan, y el stemmer puede reducirlos.
  \item Un número se pierde en la política agresiva; \textit{Ocean's 11} solo conserva \textit{ocean} y \texttt{<NUM>}.
\end{itemize}

\subsection{Casos de prueba compartidos}
El archivo \texttt{casos\_prueba/normalizacion.json} contiene 28 casos ``texto $\to$ tokens esperados'', cada uno con el resultado esperado para las tres políticas, agrupados por tema (tokenización, palabras vacías, fechas, tesauro, tildes, stemming, nombres, números, casos de borde e integración). Cualquiera de los dos indexadores debe producir exactamente esos tokens. El indexador propio los verifica con \texttt{uv run python scripts/verificar\_casos.py}; el de Lucene se verifica con la misma herramienta y la opción \texttt{--comando}, que ejecuta un programa externo por caso. El contrato completo de pasos y su orden está en \texttt{recursos/README.md}.

% =========================================================
% SECCIÓN 3: ESTADÍSTICAS
% =========================================================
\section{Estadísticas y evaluación de rendimiento de los índices}
\label{sec:estadisticas}

Las cifras de esta sección salen de scripts, no de valores escritos a mano: \texttt{scripts/estadisticas\_indices.py} lee los archivos \texttt{.idx} y genera las tablas de \texttt{tablas/} y los datos de \texttt{datos/}, que este documento incluye. Así, al reconstruir los índices solo hay que recompilar el informe. Cada índice guarda una huella de su normalización; si una política o una lista de \texttt{recursos/} cambió después de medir, el documento muestra el aviso de cifras desactualizadas.

\IfFileExists{tablas/vigencia.tex}{\newcommand{\AvisoCifras}{\par\smallskip\noindent\fbox{\parbox{0.94\linewidth}{\small\textbf{Cifras desactualizadas.} Los índices de la política minimo, estandar, avanzado se construyeron antes de los últimos cambios en el normalizador o en \texttt{recursos/}. Ejecute \texttt{scripts/regenerar\_todo.sh} y recompile.}}\par\smallskip}
\newcommand{\MedicionMeta}{; equipo no registrado}
\AvisoCifras}{}

Las tres configuraciones se construyeron sobre los mismos 500 primeros documentos de la colección, con el procedimiento de evaluación automatizado:
\begin{lstlisting}[language=bash, caption={Ejecución comparativa de las tres políticas de indexación.}]
uv run indexador benchmark --limite 500
\end{lstlisting}
La memoria de cada política se mide con la marca de agua alta del proceso (\texttt{VmHWM}), que se reinicia antes de cada ejecución; de lo contrario la segunda y la tercera política heredarían el pico de la primera.

\subsection{Métricas de los tres índices}

\begin{table}[htbp]
  \centering
  \footnotesize
  \caption{Comparativa cuantitativa de los tres índices (500 documentos).}
  \label{tab:benchmark-indices}
  \begin{tabular*}{\linewidth}{@{\extracolsep{\fill}}lrrr@{}}
    \toprule
    \textbf{Métrica} & \textbf{Mínima} & \textbf{Intermedia} & \textbf{Agresiva} \\
    \midrule
    \filas{tablas/indices.tex}
    \bottomrule
  \end{tabular*}
\end{table}

La Tabla~\ref{tab:reducciones} muestra cuánto reduce cada política frente a la mínima.

\begin{table}[htbp]
  \centering
  \footnotesize
  \caption{Reducción porcentual frente a la política mínima.}
  \label{tab:reducciones}
  \begin{tabular*}{\linewidth}{@{\extracolsep{\fill}}lrrrr@{}}
    \toprule
    \textbf{Política} & \textbf{Vocabulario} & \textbf{Postings} & \textbf{Tamaño} & \textbf{Tokens} \\
    \midrule
    \filas{tablas/reducciones.tex}
    \bottomrule
  \end{tabular*}
\end{table}

Lo que muestran las tablas:
\begin{itemize}
  \item \textbf{Las palabras vacías y el stemming dan casi toda la reducción.} Pasar de la política mínima a la intermedia es el salto grande en vocabulario, postings y tamaño; la agresiva solo añade una mejora pequeña por encima de eso.
  \item \textbf{Los tokens bajan mucho más que los postings.} Las palabras vacías son muy repetidas pero ocupan pocos postings (uno por documento), así que quitarlas reduce el total de tokens más de lo que reduce el tamaño del archivo.
  \item \textbf{El tiempo de indexación sube con la normalización.} Cada paso adicional (stemming, tesauro) cuesta tiempo por token; el índice más pequeño es el más lento de construir.
  \item \textbf{La memoria se mantiene estable} alrededor de la misma cifra en las tres políticas, porque el particionamiento de SPIMI se decide comparando la memoria residente real contra el límite del archivo de configuración y no el número de documentos.
\end{itemize}

\subsection{Distribución de los términos}

\textbf{Ley de Zipf.} La frecuencia de un término es aproximadamente inversamente proporcional a su rango~\cite{manning2008}. Se ajustó una recta a $\log f$ frente a $\log r$ en el rango \IfFileExists{tablas/zipf_rango.tex}{\newcommand{\ZipfRango}{10--5\,000}
\ZipfRango}{--}, que excluye la cabeza (dominada por palabras vacías) y la cola (dominada por términos únicos). \textbf{Ley de Heaps.} El vocabulario crece con el número de tokens como $V = K\,T^{\beta}$; los parámetros se ajustaron con la misma colección. La Tabla~\ref{tab:ajustes} resume ambos ajustes.

\begin{table}[htbp]
  \centering
  \footnotesize
  \caption{Pendiente de Zipf ($|b|$) y parámetros de Heaps ($K$, $\beta$).}
  \label{tab:ajustes}
  \begin{tabular*}{\linewidth}{@{\extracolsep{\fill}}lrrr@{}}
    \toprule
    \textbf{Política} & \textbf{Zipf $|b|$} & \textbf{Heaps $K$} & \textbf{Heaps $\beta$} \\
    \midrule
    \filas{tablas/ajustes.tex}
    \bottomrule
  \end{tabular*}
\end{table}

\IfFileExists{datos/zipf_minimo.csv}{%
\begin{figure}[htbp]
  \centering
  \begin{tikzpicture}
    \begin{loglogaxis}[width=0.9\linewidth, height=6.2cm, xlabel={Rango}, ylabel={Frecuencia}, legend pos=north east, legend style={font=\footnotesize}, grid=major, grid style={gray!25}]
      \addplot[blue, thick] table[col sep=comma, x=rango, y=frecuencia]{datos/zipf_minimo.csv};
      \addplot[red!70!black, thick, dashed] table[col sep=comma, x=rango, y=frecuencia]{datos/zipf_estandar.csv};
      \addplot[green!45!black, thick, dotted] table[col sep=comma, x=rango, y=frecuencia]{datos/zipf_avanzado.csv};
      \legend{Mínima, Intermedia, Agresiva}
    \end{loglogaxis}
  \end{tikzpicture}
  \caption{Frecuencia de los términos frente a su rango, en escala log-log.}
  \label{fig:zipf}
\end{figure}

\begin{figure}[htbp]
  \centering
  \begin{tikzpicture}
    \begin{axis}[width=0.9\linewidth, height=6.2cm, xlabel={Tokens procesados}, ylabel={Vocabulario}, legend pos=north west, legend style={font=\footnotesize}, grid=major, grid style={gray!25}, scaled ticks=false, tick label style={/pgf/number format/fixed}]
      \addplot[blue, thick] table[col sep=comma, x=tokens, y=vocabulario]{datos/heaps_minimo.csv};
      \addplot[red!70!black, thick, dashed] table[col sep=comma, x=tokens, y=vocabulario]{datos/heaps_estandar.csv};
      \addplot[green!45!black, thick, dotted] table[col sep=comma, x=tokens, y=vocabulario]{datos/heaps_avanzado.csv};
      \legend{Mínima, Intermedia, Agresiva}
    \end{axis}
  \end{tikzpicture}
  \caption{Crecimiento del vocabulario con los tokens procesados (ley de Heaps).}
  \label{fig:heaps}
\end{figure}
}{}

\textbf{Frecuencia de documento.} La Tabla~\ref{tab:topdf} lista los diez términos que aparecen en más documentos. En la política mínima son palabras vacías; en las otras dos son palabras de contenido muy generales (\textit{film}, \textit{part}, \textit{new}), y con la agresiva aparece \texttt{<NUM>}, que condensa todos los números en un solo término.

\begin{table}[htbp]
  \centering
  \footnotesize
  \caption{Diez términos con mayor frecuencia de documento (entre paréntesis, $df$).}
  \label{tab:topdf}
  \begin{tabular*}{\linewidth}{@{\extracolsep{\fill}}lP{3.4cm}P{3.4cm}P{3.4cm}@{}}
    \toprule
    \textbf{\#} & \textbf{Mínima} & \textbf{Intermedia} & \textbf{Agresiva} \\
    \midrule
    \filas{tablas/top_df.tex}
    \bottomrule
  \end{tabular*}
\end{table}

\IfFileExists{datos/df_minimo.csv}{%
\begin{figure}[htbp]
  \centering
  \begin{tikzpicture}
    \begin{axis}[ybar, width=0.95\linewidth, height=6cm, bar width=4pt, xlabel={Rango de $df$}, ylabel={Términos}, symbolic x coords={1,2--2,3--5,6--10,11--50,51--100,101--250,251--500,>500}, xtick=data, x tick label style={rotate=35, anchor=east, font=\footnotesize}, ymode=log, log origin=infty, legend style={font=\footnotesize, at={(0.98,0.97)}, anchor=north east}, enlarge x limits=0.08, ymajorgrids, grid style={gray!25}]
      \addplot[fill=blue!60] table[col sep=comma, x=rango, y=terminos]{datos/df_minimo.csv};
      \addplot[fill=red!60] table[col sep=comma, x=rango, y=terminos]{datos/df_estandar.csv};
      \addplot[fill=green!50!black] table[col sep=comma, x=rango, y=terminos]{datos/df_avanzado.csv};
      \legend{Mínima, Intermedia, Agresiva}
    \end{axis}
  \end{tikzpicture}
  \caption{Número de términos según su frecuencia de documento (escala logarítmica).}
  \label{fig:df}
\end{figure}
}{}

Cerca de la mitad del vocabulario tiene $df=1$ (Tabla~\ref{tab:benchmark-indices}). Esos términos (erratas, nombres propios raros, identificadores) casi nunca se buscan y alargan el diccionario; una política que los eliminara tendría que decidir cuáles son ruido, algo que no se hizo.

\subsection{Tamaño por sección del archivo}

La Tabla~\ref{tab:secciones} desglosa el archivo binario (estructura en la sección~\ref{sec:estructura-binaria}). Las listas de postings ocupan la mayor parte, por lo que reducir postings es lo que más achica el índice.

\begin{table}[htbp]
  \centering
  \footnotesize
  \caption{Bytes por sección del archivo \texttt{.idx}.}
  \label{tab:secciones}
  \begin{tabular*}{\linewidth}{@{\extracolsep{\fill}}lrrrrr@{}}
    \toprule
    \textbf{Política} & \textbf{Encabezado} & \textbf{Documentos} & \textbf{Diccionario} & \textbf{Postings} & \textbf{\% postings} \\
    \midrule
    \filas{tablas/secciones.tex}
    \bottomrule
  \end{tabular*}
\end{table}

\subsection{Consultas bajo cada política}

La Tabla~\ref{tab:consultas} cuenta los documentos que devuelve cada política para consultas de tipos distintos, y la Tabla~\ref{tab:consultas-tokens} muestra en qué términos se convirtió cada una.

\begin{table}[htbp]
  \centering
  \footnotesize
  \caption{Documentos que contienen al menos un término de la consulta.}
  \label{tab:consultas}
  \begin{tabular*}{\linewidth}{@{\extracolsep{\fill}}lrrr@{}}
    \toprule
    \textbf{Consulta} & \textbf{Mínima} & \textbf{Intermedia} & \textbf{Agresiva} \\
    \midrule
    \filas{tablas/consultas.tex}
    \bottomrule
  \end{tabular*}
\end{table}

\begin{table}[htbp]
  \centering
  \footnotesize
  \caption{Términos en que se convierte cada consulta.}
  \label{tab:consultas-tokens}
  \begin{tabular*}{\linewidth}{@{\extracolsep{\fill}}lP{3.3cm}P{3.3cm}P{3.3cm}@{}}
    \toprule
    \textbf{Consulta} & \textbf{Mínima} & \textbf{Intermedia} & \textbf{Agresiva} \\
    \midrule
    \filas{tablas/consultas_tokens.tex}
    \bottomrule
  \end{tabular*}
\end{table}

Estas cifras cuentan documentos que contienen \emph{algún} término de la consulta, no resultados relevantes: miden cobertura, no precisión, y no hay juicios de relevancia con los que medir esta última. Aun así se ven efectos claros. Con la agresiva, \textit{película} y \textit{movie} se unifican en \texttt{film} y devuelven el mismo conjunto en ambos idiomas; \textit{director} y \textit{cineasta} también convergen. En cambio \textit{Star Wars} devuelve muchos más documentos con la intermedia que con la mínima porque el stemmer une \textit{star} y \textit{stars}, y la consulta de un año (\textit{1994}) se vuelve \texttt{<ANIO\_1994>}, pero con la agresiva baja de 38 a 27 documentos: un año que aparece dentro de una fecha completa se reemplaza junto con ella por \texttt{<FECHA>} y deja de ser buscable.

\subsection{Comparación entre el indexador propio y el de Lucene}

El script \texttt{scripts/benchmark\_comparativo.py} ejecuta ambos indexadores como procesos externos con las mismas tres configuraciones y la misma muestra, y mide cada uno con \texttt{RUSAGE\_CHILDREN} desde un proceso intermedio. Así Python y la JVM se miden de la misma manera: tiempo de pared, tiempo de CPU, memoria residente máxima, tamaño del índice, términos y postings.

\IfFileExists{tablas/comparativo.tex}{%
\begin{table}[htbp]
  \centering
  \footnotesize
  \caption{Indexador propio frente a Lucene.}
  \label{tab:comparativo}
  \begin{tabular*}{\linewidth}{@{\extracolsep{\fill}}llrrrrrr@{}}
    \toprule
    \textbf{Impl.} & \textbf{Política} & \textbf{Pared (s)} & \textbf{CPU (s)} & \textbf{RAM (MB)} & \textbf{Tam. (MB)} & \textbf{Términos} & \textbf{Postings} \\
    \midrule
    \filas{tablas/comparativo.tex}
    \bottomrule
  \end{tabular*}
\end{table}
}{%
\noindent\fbox{\parbox{0.94\linewidth}{\small \textbf{Pendiente.} La comparación con Lucene aún no tiene cifras: falta el indexador de Lucene. Cuando exista, se mide con \texttt{scripts/regenerar\_todo.sh} definiendo la variable \texttt{LUCENE\_CMD}, y esta tabla aparece sola.}}
}

\subsection{Composición de la muestra y razón texto-índice}

\IfFileExists{tablas/dominios.tex}{%
La muestra de 500 documentos reparte su texto de forma muy desigual entre los dominios de origen (Tabla~\ref{tab:dominios}). Los guiones de \texttt{dailyscript.com} son solo el $16{,}4\,\%$ de los documentos, pero aportan más de la mitad del texto, mientras que \texttt{filmaffinity.com} y \texttt{tvmaze.com} contribuyen con fichas muy cortas.

\begin{table}[htbp]
  \centering
  \footnotesize
  \caption{Documentos y texto de la muestra por dominio de origen.}
  \label{tab:dominios}
  \begin{tabular*}{\linewidth}{@{\extracolsep{\fill}}lrrr@{}}
    \toprule
    \textbf{Dominio} & \textbf{Documentos} & \textbf{\% docs.} & \textbf{Texto (MB)} \\
    \midrule
    \filas{tablas/dominios.tex}
    \bottomrule
  \end{tabular*}
\end{table}
}{}

\IfFileExists{tablas/razon_texto.tex}{%
La Tabla~\ref{tab:razon-texto} compara el tamaño de cada índice con el texto del que sale. Incluso la política mínima, que conserva todas las palabras, ocupa poco más de la mitad del texto original, porque el índice guarda cada término una sola vez y solo un registro de doce bytes por par término-documento. Las políticas intermedia y agresiva lo bajan a cerca del $41\,\%$ y el $40\,\%$.

\begin{table}[htbp]
  \centering
  \footnotesize
  \caption{Tamaño del índice frente al texto de la muestra.}
  \label{tab:razon-texto}
  \begin{tabular*}{\linewidth}{@{\extracolsep{\fill}}lrrr@{}}
    \toprule
    \textbf{Política} & \textbf{Texto (MB)} & \textbf{Índice (MB)} & \textbf{Índice / texto (\%)} \\
    \midrule
    \filas{tablas/razon_texto.tex}
    \bottomrule
  \end{tabular*}
\end{table}
}{}

\subsection{Escala de la colección}

Una prueba de escalado sobre veinte mil documentos de la misma colección confirmó que el esquema se sostiene fuera de la muestra: el índice se construyó en $160{,}8$ segundos con $539\,511$ términos y $13\,637\,739$ ocurrencias, un archivo de $172{,}6$ MB y un pico de $552$ MB de memoria residente, lo que proyecta alrededor de una hora y $3{,}7$ GB para los $317\,657$ documentos de la colección ($5{,}10$ GiB de texto). Esa medición motiva que el límite de memoria de SPIMI se evalúe contra el consumo real del proceso y no contra un número fijo de documentos.


% =========================================================
% SECCIÓN 4: ESTRUCTURA INTERNA DEL ARCHIVO BINARIO
% =========================================================
\section{Descripción de la estructura interna del archivo del índice}
\label{sec:estructura-binaria}

\subsection{Fundamentación del formato binario manual}

En correspondencia con las directrices de la especificación técnica, el índice desarrollado prescinde de bases de datos embebidas o módulos de serialización automática, administrándose como un archivo binario monolítico con extensión \texttt{.idx} que se escribe y lee directamente a nivel de bytes mediante empaquetamiento Little-Endian con el módulo \texttt{struct}.

La arquitectura del formato se fundamenta en dos principios rectores:
\begin{enumerate}
    \item \textbf{Alineación deliberada para inspección con \texttt{hexdump -C}:} el encabezado se reservó en sesenta y cuatro bytes exactos para cubrir con precisión las primeras cuatro líneas de dieciséis bytes del visor hexadecimal, mientras que los registros documentales y de ocurrencias poseen dimensiones fijas de cuarenta y cuatro y doce bytes respectivamente, almacenando las cadenas del vocabulario en codificación UTF-8 para su lectura directa en la columna ASCII.
    \item \textbf{Capacidad de acceso aleatorio:} la uniformidad en las estructuras permite consultar la norma vectorial o los metadatos de cualquier documento en tiempo constante calculando su posición física como $\text{offset} = 64 + (\text{doc\_id} \times 44)$, al tiempo que la ordenación alfabética del léxico posibilita búsquedas binarias en $O(\log |V|)$ y saltos directos hacia los bloques de ocurrencias asociados.
\end{enumerate}

\begin{figure}[htbp]
  \centering
  \begin{tikzpicture}[
    scale=0.96,
    every node/.style={transform shape},
    x=1cm, y=1cm,
    font=\small,
    secnode/.style={
      rectangle, draw=black!70, line width=0.6pt, rounded corners=3pt,
      inner sep=0pt
    },
    cell/.style={
      rectangle, draw=black!55, line width=0.5pt, fill=white, rounded corners=1.5pt,
      font=\scriptsize, align=center, inner sep=2pt, minimum height=0.72cm
    },
    offsetlabel/.style={
      font=\ttfamily\scriptsize\bfseries\color{blue!80!black},
      anchor=east
    },
    secnum/.style={
      font=\bfseries\footnotesize,
      anchor=south west
    }
  ]
    % -----------------------------------------------------------
    % Sección 1: Encabezado global (64 bytes)
    % -----------------------------------------------------------
    \node[secnum, color=blue!75!black] at (0, 1.53) {Sección 1};
    \node[secnode, fill=blue!8, minimum width=12.6cm, minimum height=3.0cm] (s1) at (6.3, 0) {};
    \node[anchor=west, font=\bfseries\footnotesize] at (0.35, 1.1) {Encabezado global de 64 bytes};
    \node[anchor=east, font=\scriptsize\color{black!60}] at (12.25, 1.1) {0x00 a 0x3F};
    
    % Fila 1 de la cabecera (28 bytes): ancho proporcional a 8+2+2+4+4+8
    \node[cell, minimum width=3.371cm] at (1.7855, 0.35) {Firma \texttt{ENCUADRE}\\ \tiny\color{black!70} 8 bytes};
    \node[cell, minimum width=0.843cm, font=\tiny] at (4.0125, 0.35) {Ver.\\ \tiny\color{black!70} 2 bytes};
    \node[cell, minimum width=0.843cm, font=\tiny] at (4.9755, 0.35) {Flags\\ \tiny\color{black!70} 2 bytes};
    \node[cell, minimum width=1.686cm] at (6.36, 0.35) {Docs $N$\\ \tiny\color{black!70} 4 bytes};
    \node[cell, minimum width=1.686cm] at (8.166, 0.35) {Léxico $|V|$\\ \tiny\color{black!70} 4 bytes};
    \node[cell, minimum width=3.371cm] at (10.8145, 0.35) {Postings\\ \tiny\color{black!70} 8 bytes};

    % Fila 2 de la cabecera (36 bytes): tres offsets y el modelo de ponderación
    \node[cell, minimum width=2.622cm] at (1.411, -0.65) {Offset documentos\\ \tiny\color{black!70} 8 bytes · \texttt{0x00000040}};
    \node[cell, minimum width=2.622cm] at (4.153, -0.65) {Offset diccionario\\ \tiny\color{black!70} 8 bytes · \texttt{0x00005630}};
    \node[cell, minimum width=2.622cm] at (6.895, -0.65) {Offset postings\\ \tiny\color{black!70} 8 bytes · \texttt{0x001B2423}};
    \node[cell, minimum width=3.933cm] at (10.2925, -0.65) {Modelo $tf$/$idf$ y banderas\\ \tiny\color{black!70} 12 bytes · \texttt{0x00}};

    % -----------------------------------------------------------
    % Sección 2: Tabla de documentos (500 docs x 44 bytes)
    % -----------------------------------------------------------
    \node[secnum, color=green!55!black] at (0, -2.07) {Sección 2};
    \node[secnode, fill=green!8, minimum width=12.6cm, minimum height=2.1cm] (s2) at (6.3, -3.15) {};
    \node[anchor=west, font=\bfseries\footnotesize] at (0.35, -2.5) {Tabla de documentos y normas vectoriales};
    \node[anchor=east, font=\scriptsize\color{black!60}] at (12.25, -2.5) {500 registros · 44 bytes cada uno};
    
    % 4 celdas con separación fija de 2 mm
    \node[cell, minimum width=2.20cm] at (1.40, -3.45) {\texttt{doc\_id}\\ \tiny\color{black!70} 4 bytes};
    \node[cell, minimum width=4.60cm] at (5.00, -3.45) {\texttt{hash\_url} · SHA-256 binario\\ \tiny\color{black!70} 32 bytes};
    \node[cell, minimum width=2.30cm] at (8.65, -3.45) {\texttt{num\_tokens}\\ \tiny\color{black!70} 4 bytes};
    \node[cell, minimum width=2.30cm] at (11.15, -3.45) {Norma $\|d\|_2$\\ \tiny\color{black!70} 4 bytes · IEEE 754};

    % -----------------------------------------------------------
    % Sección 3: Diccionario (59 224 términos)
    % -----------------------------------------------------------
    \node[secnum, color=orange!70!black] at (0, -4.77) {Sección 3};
    \node[secnode, fill=orange!8, minimum width=12.6cm, minimum height=2.1cm] (s3) at (6.3, -5.85) {};
    \node[anchor=west, font=\bfseries\footnotesize] at (0.35, -5.2) {Diccionario de términos};
    \node[anchor=east, font=\scriptsize\color{black!60}] at (12.25, -5.2) {$60\,672$ términos en orden alfabético};
    
    % 6 celdas con separación fija de 2 mm
    \node[cell, minimum width=1.70cm] at (1.15, -6.15) {\texttt{len\_termino}\\ \tiny\color{black!70} 2 bytes};
    \node[cell, minimum width=3.10cm] at (3.75, -6.15) {Término UTF-8\\ \tiny\color{black!70} longitud variable};
    \node[cell, minimum width=1.40cm] at (6.20, -6.15) {$df$\\ \tiny\color{black!70} 4 bytes};
    \node[cell, minimum width=1.40cm] at (7.80, -6.15) {$IDF$\\ \tiny\color{black!70} 4 bytes};
    \node[cell, minimum width=2.10cm] at (9.75, -6.15) {\texttt{offset\_postings}\\ \tiny\color{black!70} 8 bytes};
    \node[cell, minimum width=1.30cm] at (11.65, -6.15) {\texttt{bytes}\\ \tiny\color{black!70} 4 bytes};

    % -----------------------------------------------------------
    % Sección 4: Postings (459 396 tuplas x 12 bytes)
    % -----------------------------------------------------------
    \node[secnum, color=purple!70!black] at (0, -7.47) {Sección 4};
    \node[secnode, fill=purple!8, minimum width=12.6cm, minimum height=2.1cm] (s4) at (6.3, -8.55) {};
    \node[anchor=west, font=\bfseries\footnotesize] at (0.35, -7.9) {Listas contiguas de ocurrencias};
    \node[anchor=east, font=\scriptsize\color{black!60}] at (12.25, -7.9) {$479\,680$ tuplas · 12 bytes cada una};
    
    % 3 celdas con separación fija de 2 mm
    \node[cell, minimum width=3.30cm] at (1.95, -8.85) {\texttt{doc\_id}\\ \tiny\color{black!70} 4 bytes · uint32};
    \node[cell, minimum width=3.30cm] at (5.45, -8.85) {$tf$ · frecuencia bruta\\ \tiny\color{black!70} 4 bytes · uint32};
    \node[cell, minimum width=5.00cm] at (9.80, -8.85) {Ponderación $w_{t,d} = (1 + \ln(tf)) \times IDF$\\ \tiny\color{black!70} 4 bytes · IEEE 754};

    % -----------------------------------------------------------
    % Eje vertical de desplazamientos en disco
    % -----------------------------------------------------------
    \draw[line width=0.8pt, blue!70!black] (-0.35, 1.5) -- (-0.35, -9.6);

    \foreach \y/\txt in {
      1.5/{0x00000000},
      -2.1/{0x00000040},
      -4.8/{0x00005630},
      -7.5/{0x001B2423},
      -9.6/{0x0072F923}
    } {
      \draw[line width=0.8pt, blue!70!black] (-0.50, \y) -- (-0.20, \y);
      \node[offsetlabel] at (-0.60, \y) {\txt};
    }
  \end{tikzpicture}
  \caption{Disposición secuencial de las cuatro secciones en el archivo binario \texttt{.idx}.}
  \label{fig:layout-binario}
\end{figure}

\subsection{Encabezado global de 64 bytes}

Los primeros sesenta y cuatro bytes del archivo corresponden a los metadatos globales del índice, situados en el intervalo comprendido entre \texttt{0x00000000} y \texttt{0x0000003F}, accesibles mediante el comando \texttt{hexdump -C -n 64 salida/indice\_estandar.idx}:

\begin{lstlisting}[language=bash, caption={Volcado del encabezado mediante \texttt{hexdump}.}]
00000000  45 4e 43 55 41 44 52 45  01 00 0f 00 f4 01 00 00  |ENCUADRE........|
00000010  00 ed 00 00 c0 51 07 00  00 00 00 00 40 00 00 00  |.....Q......@...|
00000020  00 00 00 00 30 56 00 00  00 00 00 00 23 24 1b 00  |....0V......#$..|
00000030  00 00 00 00 00 00 00 00  00 00 00 00 00 00 00 00  |................|
\end{lstlisting}

\begin{table}[htbp]
  \centering
  \small
  \caption{Mapeo campo por campo del encabezado del archivo \texttt{.idx}.}
  \label{tab:encabezado-binario}
  \begin{tabularx}{\linewidth}{llrY}
    \toprule
    \textbf{Desplazamiento} & \textbf{Tipo} & \textbf{Bytes} & \textbf{Interpretación y valor decodificado} \\
    \midrule
    \texttt{0x00 - 0x07} & \texttt{char[8]} & 8 & Firma mágica de identificación: \texttt{ENCUADRE} \\
    \texttt{0x08 - 0x09} & \texttt{uint16}  & 2 & Versión del esquema binario: 1 \\
    \texttt{0x0A - 0x0B} & \texttt{uint16}  & 2 & Banderas de normalización: 0x000F con minúsculas, puntuación, palabras vacías y stemming Snowball activados \\
    \texttt{0x0C - 0x0F} & \texttt{uint32}  & 4 & Total de documentos $N$: 500 \\
    \texttt{0x10 - 0x13} & \texttt{uint32}  & 4 & Tamaño del vocabulario $|V|$: $60\,672$ términos \\
    \texttt{0x14 - 0x1B} & \texttt{uint64}  & 8 & Total de ocurrencias acumuladas: $479\,680$ postings \\
    \texttt{0x1C - 0x23} & \texttt{uint64}  & 8 & Desplazamiento a la tabla de documentos: byte 64, equivalente a 0x40 \\
    \texttt{0x24 - 0x2B} & \texttt{uint64}  & 8 & Desplazamiento a la sección del diccionario: byte $22\,064$, equivalente a 0x5630 \\
    \texttt{0x2C - 0x33} & \texttt{uint64}  & 8 & Desplazamiento a la sección de postings: byte $1\,778\,723$, equivalente a 0x1B2423 \\
    \texttt{0x34 - 0x3F} & \texttt{byte[12]} & 12 & Modelo de ponderación: byte 0 con el esquema $tf$, byte 1 con el esquema $idf$ y byte 2 con banderas de coseno e IDF suavizado; los nueve bytes restantes quedan en cero \\
    \bottomrule
  \end{tabularx}
\end{table}

\subsection{Tabla de documentos y normas vectoriales}

A partir de la posición \texttt{0x00000040} se despliega la tabla continua de documentos y normas vectoriales, conformada por entradas contiguas de cuarenta y cuatro bytes empaquetadas bajo el formato \texttt{<I32sIf}:
\begin{itemize}
    \item \texttt{doc\_id}: entero sin signo de 32 bits y 4 bytes que representa el identificador secuencial continuo.
    \item \texttt{hash\_url}: arreglo de 32 bytes con la clave SHA-256 binaria del documento.
    \item \texttt{num\_tokens}: entero sin signo de 32 bits y 4 bytes con el conteo total de tokens del documento.
    \item \texttt{norma\_vectorial}: flotante de 4 bytes con la norma euclidiana $\|d\|_2 = \sqrt{\sum w_{t,d}^2}$.
\end{itemize}

\begin{lstlisting}[language=bash, caption={Volcado de los dos primeros registros documentales mediante \texttt{hexdump}.}]
00000040  00 00 00 00 df 3c 0a bb  ab e9 ad 1e 5a a2 0b eb  |.....<......Z...|
00000050  3e f1 5f 63 59 07 48 16  c4 71 75 14 6d fa a8 20  |>._cY.H..qu.m.. |
00000060  ac ce 46 b6 e3 34 00 00  4c f3 9d 43 01 00 00 00  |..F..4..L..C....|
00000070  08 67 35 fc f8 5c 56 b5  f3 fd 10 2c 83 71 19 7a  |.g5..\V....,.q.z|
00000080  bf a2 11 ec  21 90 9d ff  5a b5 99 da 83 37 0c c2  |....!...Z....7..|
00000090  5b 0d 00 00 e9 13 32 43                           |.[.....2C|
\end{lstlisting}

Al examinar detalladamente los primeros dos registros del volcado anterior, se aprecian sus componentes correspondientes:
\begin{itemize}
    \item \textbf{Documento 0 en el rango \texttt{0x00000040} a \texttt{0x0000006B}:}
    identificador secuencial cero, resumen criptográfico SHA-256 almacenado en treinta y dos bytes, longitud calculada de $13\,539$ tokens y una norma vectorial $\|d\|_2 = 315.9008$ representada por la secuencia en coma flotante \texttt{4c f3 9d 43}.
    \item \textbf{Documento 1 en el rango \texttt{0x0000006C} a \texttt{0x00000097}:}
    identificador secuencial uno, resumen SHA-256, longitud de $3\,419$ tokens y norma $\|d\|_2 = 178.0778$ codificada en los bytes \texttt{e9 13 32 43}.
\end{itemize}

\subsection{Diccionario o vocabulario}

En la dirección \texttt{0x00005630}, obtenida al sumar el bloque de cabecera y las quinientas entradas documentales previas, comienza la sección del léxico, donde cada término se estructura con una longitud dinámica según la extensión de su cadena:
\begin{itemize}
    \item \texttt{len\_termino}: entero de 16 bits y 2 bytes con la longitud de la cadena en UTF-8.
    \item \texttt{termino\_utf8}: bytes legibles directamente en la columna ASCII del visor hexadecimal.
    \item \texttt{df}: entero de 32 bits y 4 bytes con la frecuencia documental.
    \item \texttt{idf}: flotante de 4 bytes con el peso inverso documental calculado según $\ln(N/df + 1)$.
    \item \texttt{offset\_postings}: entero de 64 bits y 8 bytes con el desplazamiento absoluto al bloque de ocurrencias.
    \item \texttt{bytes\_postings}: entero de 32 bits y 4 bytes con la longitud en bytes de la lista.
\end{itemize}

\begin{lstlisting}[language=bash, caption={Volcado del diccionario de términos mediante \texttt{hexdump}.}]
00005630  03 00 30 2d 30 01 00 00  00 70 ee c6 40 23 24 1b  |..0-0....p..@#$.|
00005640  00 00 00 00 00 0c 00 00  00 0d 00 30 2d 31 32 2d  |...........0-12-|
00005650  34 39 30 36 33 30 2d 33  01 00 00 00 70 ee c6 40  |490630-3....p..@|
00005660  2f 24 1b 00 00 00 00 00  0c 00 00 00 03 00 30 2d  |/$............0-|
00005670  32 01 00 00 00 70 ee c6  40 3b 24 1b 00 00 00 00  |2....p..@;$.....|
\end{lstlisting}

Al contrastar las entradas consecutivas capturadas en el volcado anterior, se comprueba la serialización individual de cada término:
\begin{itemize}
    \item \textbf{Término \texttt{"0-0"}:} longitud de 3 bytes con valor \texttt{03 00}, cadena ASCII legible \texttt{"0-0"}, $df = 1$, peso inverso $IDF = \ln(500/1 + 1) = 6.2166$ en los bytes \texttt{70 ee c6 40}, desplazamiento de postings en \texttt{0x001B2423} y tamaño de postings de 12 bytes.
    \item \textbf{Término \texttt{"0-12-490630-3"}:} longitud de 13 bytes con valor \texttt{0d 00}, caracteres visibles en el visor ASCII \texttt{"0-12-490630-3"} y desplazamiento de ocurrencias en \texttt{0x001B242F}.
    \item \textbf{Término \texttt{"0-2"}:} longitud de 3 bytes con valor \texttt{03 00}, carácter visible \texttt{"0-2"} y desplazamiento de ocurrencias en \texttt{0x001B243B}.
\end{itemize}

\subsection{Listas de ocurrencias o postings}

A partir del byte \texttt{0x001B2423} se despliegan las listas de ocurrencias agrupadas por término, en las cuales cada posting individual se compacta en doce bytes según el patrón \texttt{<IIf}:
\begin{itemize}
    \item \texttt{doc\_id}: identificador de documento de 32 bits y 4 bytes.
    \item \texttt{tf}: frecuencia bruta del término en el documento, en 4 bytes.
    \item \texttt{peso\_tfidf}: peso ponderado $w_{t,d} = (1 + \ln(tf)) \times IDF$ en 4 bytes.
\end{itemize}

\begin{lstlisting}[language=bash, caption={Volcado de la lista de ocurrencias mediante \texttt{hexdump}.}]
001b2423  1b 01 00 00 02 00 00 00  fa 68 28 41 8a 00 00 00  |.........h(A....|
001b2433  01 00 00 00 70 ee c6 40  4d 00 00 00 04 00 00 00  |....p..@M.......|
001b2443  bc 5a 6d 41 85 01 00 00  01 00 00 00 70 ee c6 40  |.ZmA........p..@|
\end{lstlisting}

\begin{table}[htbp]
  \centering
  \small
  \caption{Mapeo de registros de postings consecutivos.}
  \label{tab:postings-binarios}
  \begin{tabular*}{\linewidth}{@{\extracolsep{\fill}}lrrllr@{}}
    \toprule
    \textbf{Offset hex} & \textbf{Doc} & \textbf{TF} & \textbf{Hex IEEE 754} & \textbf{Peso $w_{t,d}$} & \textbf{Término} \\
    \midrule
    \texttt{0x001B2423} & \texttt{283} & 2 & \texttt{fa 68 28 41} & $10.5256$ & \texttt{"0-0"} \\
    \texttt{0x001B242F} & \texttt{138} & 1 & \texttt{70 ee c6 40} & $6.2166$ & \texttt{"0-12-490630-3"} \\
    \texttt{0x001B243B} & \texttt{77}  & 4 & \texttt{bc 5a 6d 41} & $14.8347$ & \texttt{"0-2"} \\
    \texttt{0x001B2447} & \texttt{389} & 1 & \texttt{70 ee c6 40} & $6.2166$ & \texttt{"0-34"} \\
    \bottomrule
  \end{tabular*}
\end{table}

La correspondencia del modelo matemático se comprueba verificando la coincidencia entre la ponderación teórica y su registro binario; por ejemplo, para el término \texttt{"0-0"} presente en el documento doscientos ochenta y tres con una frecuencia $tf=2$:
\begin{equation}
    w_{t,d} = (1 + \ln(2)) \times 6.2166 = (1 + 0.693147) \times 6.2166 \approx 10.5256
\end{equation}
cuyo valor en coma flotante Little-Endian IEEE 754 genera con exactitud la secuencia \texttt{fa 68 28 41}, equivalente a \texttt{0x412868fa} y presente de forma idéntica en el archivo físico.

\subsection{Inspección automatizada de la estructura binaria}

Complementariamente a la lectura directa de bytes mediante \texttt{hexdump}, el indexador provee una herramienta de línea de comandos diseñada para verificar la consistencia del archivo, comprobar los desplazamientos calculados y recuperar los postings de cualquier término en tiempo $O(\log |V|)$ sin cargar el índice completo a memoria:
\begin{lstlisting}[language=bash, caption={Inspección de estructuras y términos en el archivo binario.}]
uv run indexador inspeccionar salida/indice_estandar.idx --termino "film"
\end{lstlisting}


% =========================================================
% REFERENCIAS
% =========================================================
\begin{thebibliography}{99}
\setlength{\itemsep}{4pt}
\bibitem{manning2008} C.~D. Manning, P.~Raghavan y H.~Schütze, \textit{Introduction to Information Retrieval}. Cambridge University Press, 2008. Secciones 2.2 (normalización y stemming), 5.1 (leyes de Heaps y Zipf) y 6.2 (ponderación). \url{https://nlp.stanford.edu/IR-book/}
\bibitem{porter1980} M.~F. Porter, ``An algorithm for suffix stripping'', \textit{Program}, vol.~14, n.º~3, pp.~130--137, 1980.
\bibitem{porter2001} M.~F. Porter, ``Snowball: A language for stemming algorithms'', octubre de 2001. \url{https://snowballstem.org/texts/introduction.html}
\bibitem{harman1991} D.~Harman, ``How effective is suffixing?'', \textit{Journal of the American Society for Information Science}, vol.~42, n.º~1, pp.~7--15, 1991. \url{https://doi.org/10.1002/(SICI)1097-4571(199101)42:1<7::AID-ASI2>3.0.CO;2-P}
\bibitem{voorhees1994} E.~M. Voorhees, ``Query expansion using lexical-semantic relations'', en \textit{Proc.\ 17th ACM SIGIR}, Dublín, pp.~61--69, 1994.
\bibitem{salton1988} G.~Salton y C.~Buckley, ``Term-weighting approaches in automatic text retrieval'', \textit{Information Processing \& Management}, vol.~24, n.º~5, pp.~513--523, 1988. \url{https://doi.org/10.1016/0306-4573(88)90021-0}
\bibitem{unicode15} Unicode Consortium, ``Unicode Standard Annex \#15: Unicode Normalization Forms''. \url{https://unicode.org/reports/tr15/}
\bibitem{iso25964} ISO, \textit{ISO 25964-1:2011, Information and documentation -- Thesauri and interoperability with other vocabularies -- Part 1: Thesauri for information retrieval}, 2011. \url{https://www.iso.org/standard/53657.html}
\bibitem{nadeau2007} D.~Nadeau y S.~Sekine, ``A survey of named entity recognition and classification'', \textit{Lingvisticae Investigationes}, vol.~30, n.º~1, pp.~3--26, 2007. \url{https://doi.org/10.1075/li.30.1.03nad}
\bibitem{lucene-wordlist} Apache Lucene, documentación de la clase \texttt{WordlistLoader}. \url{https://lucene.apache.org/core/9_10_0/core/org/apache/lucene/analysis/WordlistLoader.html}
\bibitem{lucene-solrsyn} Apache Lucene, documentación de la clase \texttt{SolrSynonymParser}. \url{https://lucene.apache.org/core/8_1_0/analyzers-common/org/apache/lucene/analysis/synonym/SolrSynonymParser.html}
\end{thebibliography}

\end{document}
